\begin{question}
\noindent A telescope manufacturer produces two similar cone-shaped mirror blanks, a small prototype and a full-size mirror, for testing before grinding into parabolic shape.

\begin{center}
\begin{tikzpicture}[scale=0.7]
    % small cone
    \draw (0,0) -- (1.6,0) arc (0:-180:0.8 and 0.25) -- cycle;
    \draw (0.8,0.25) -- (0.8,3.2);
    \draw (0,0) -- (0.8,3.2) -- (1.6,0);
    \draw[dashed] (0,0) arc (180:360:0.8 and 0.25);
    \node at (0.8,-0.7) {8 cm};
    \node at (2.3,1.6) {small mirror};

    % large cone
    \draw (5,0) -- (9,0) arc (0:-180:2 and 0.5) -- cycle;
    \draw (7,0.5) -- (7,8);
    \draw (5,0) -- (7,8) -- (9,0);
    \draw[dashed] (5,0) arc (180:360:2 and 0.5);
    \node at (7,-0.9) {20 cm};
    \node at (9.8,4.5) {full-size mirror};
\end{tikzpicture}
\end{center}

\noindent The small mirror blank has a base diameter of 8 cm. The full-size mirror blank has a base diameter of 20 cm. The two blanks are mathematically similar cones.

\noindent The curved surface area of the small mirror blank is $150.8 \text{ cm}^2$.
\noindent The volume of the full-size mirror blank is $2400 \text{ cm}^3$.

\begin{enumerate}
    \item[(a)] Find the linear scale factor from the small mirror blank to the full-size mirror blank.

    \vspace{2.5cm}
    \answerplain{1}

    \item[(b)] Calculate the curved surface area of the full-size mirror blank. Give your answer to 3 significant figures.

    \vspace{4cm}
    \answerunit{$\text{cm}^2$}{2}

    \item[(c)] Calculate the volume of the small mirror blank. Give your answer to 3 significant figures.

    \vspace{4cm}
    \answerunit{$\text{cm}^3$}{2}
\end{enumerate}
\end{question}
